% !TEX root = ../main.tex

\section{Transformers and attention}\label{sec:attention}

\begin{frame}{From token vectors to contextual vectors}
  \begin{center}
    \begin{tikzpicture}[
      stage/.style={draw, rounded corners=3pt, minimum width=2.1cm,
        minimum height=0.8cm, align=center},
      >=Latex
    ]
      \node[stage, fill=tokenblue, draw=tokenblueborder] (tokens) {tokens};
      \node[stage, fill=tokenorange, draw=tokenorangeborder,
        right=0.55cm of tokens] (ids) {IDs};
      \node[stage, fill=tokengreen, draw=tokengreenborder,
        right=0.55cm of ids] (vectors) {embeddings};
      \node[stage, fill=tensorpurple, draw=tensorpurpleborder,
        right=0.55cm of vectors] (transformer) {Transformer\\?};
      \draw[->, thick] (tokens) -- (ids);
      \draw[->, thick] (ids) -- (vectors);
      \draw[->, thick] (vectors) -- (transformer);
    \end{tikzpicture}
  \end{center}

  \vspace{0.8em}
  The embedding lookup gives \tokenbox{bank} the same starting vector in both
  sentences:
  \begin{itemize}
    \item ``They sat on the river \textbf{bank}''.
    \item ``She deposited money at the \textbf{bank}''.
  \end{itemize}

  \begin{alertblock}{The problem self-attention solves}
    How can every token exchange information with the other tokens so that its
    vector reflects the surrounding context? The following slides build the
    mathematical mechanism that makes this possible.
  \end{alertblock}
\end{frame}

\begin{frame}{Attention creates a context-dependent update
  \hfill{\tiny \texttt{Head.forward}: lines 15--31}}
  \small
  The lookup vector for \texttt{bank} starts identically in both sentences.
  Attention lets it collect different clues from the surrounding tokens; we
  denote the resulting contextual vector by $\mathbf h$.

  \begin{columns}[T]
    \begin{column}{0.48\textwidth}
      \begin{block}{River context}
        \centering
        \begin{tikzpicture}[>=Latex, node distance=0.18cm]
          \node (river) {\tokenbox{river}};
          \node[right=of river] (bank) {\tokenbox[tokengreen]{bank}};
          \draw[->, very thick, tokenblueborder]
            (bank.south) .. controls +(-0.35,-0.12) and +(0.35,-0.12) ..
            (river.south);
        \end{tikzpicture}

        \[
          \mathbf e_{\text{bank}}
          \longrightarrow
          \mathbf h_{\text{bank}}^{(\text{river})}
        \]
        The result incorporates information associated with a river edge.
      \end{block}
    \end{column}
    \begin{column}{0.48\textwidth}
      \begin{block}{Financial context}
        \centering
        \begin{tikzpicture}[>=Latex, node distance=0.18cm]
          \node (money) {\tokenbox{money}};
          \node[right=of money] (bank) {\tokenbox[tokengreen]{bank}};
          \draw[->, very thick, tokenorangeborder]
            (bank.south) .. controls +(-0.35,-0.12) and +(0.35,-0.12) ..
            (money.south);
        \end{tikzpicture}

        \[
          \mathbf e_{\text{bank}}
          \longrightarrow
          \mathbf h_{\text{bank}}^{(\text{money})}
        \]
        The result incorporates information associated with finance.
      \end{block}
    \end{column}
  \end{columns}

  \vspace{-0.4em}
  \begin{alertblock}{The central idea}
    \footnotesize
    Each position asks which other positions matter, then mixes information
    from them. The resulting vector is contextual rather than a static lookup.
  \end{alertblock}
\end{frame}

\begin{frame}{Before attention: the model needs token order
  \hfill{\tiny \texttt{BigramLanguageModel.forward}: lines 18--23}}
  \small
  A sentence is not merely a bag of tokens: changing the order changes the
  meaning. A basic construction adds a position vector to every token vector:
  \[
    \mathbf x_i=\mathbf e_{\text{token}_i}+\mathbf p_i,
    \qquad X\in\mathbb R^{n\times d}.
  \]

  \begin{center}
    \begin{tikzpicture}[
      box/.style={draw, rounded corners=3pt, minimum width=2.1cm,
        minimum height=0.75cm, align=center},
      >=Latex
    ]
      \node[box, fill=tokenblue, draw=tokenblueborder] (token)
        {token vector\\$\mathbf e_i$};
      \node[box, fill=tokenorange, draw=tokenorangeborder,
        right=1.0cm of token] (position) {position vector\\$\mathbf p_i$};
      \node[right=0.42cm of token] (plus) {$+$};
      \node[box, fill=tokengreen, draw=tokengreenborder,
        right=1.0cm of position] (input) {input row\\$\mathbf x_i$};
      \draw[->, thick] (position) -- (input);
    \end{tikzpicture}
  \end{center}

  \begin{block}{Implementations differ}
    Position information may be learned or fixed. Rotary position encodings
    instead inject relative position when queries and keys are compared. The
    shared goal is the same: make order available to attention.
  \end{block}
\end{frame}

\begin{frame}{Query, key, value: three roles for each token
  \hfill{\tiny \texttt{Head}: lines 8--10, 17--30}}
  \small
  Self-attention derives three vectors from every input position.

  \begin{center}
    \begin{tikzpicture}[
      role/.style={draw, rounded corners=4pt, minimum width=3.35cm,
        minimum height=1.45cm, align=center},
      >=Latex
    ]
      \node[role, fill=tokenorange, draw=tokenorangeborder] (query)
        {\textbf{Query $\mathbf q_i$}\\What information do\\I need?};
      \node[role, fill=tokenblue, draw=tokenblueborder,
        right=0.45cm of query] (key)
        {\textbf{Key $\mathbf k_j$}\\What kind of information\\do I offer?};
      \node[role, fill=tokengreen, draw=tokengreenborder,
        right=0.45cm of key] (value)
        {\textbf{Value $\mathbf v_j$}\\What information will\\I contribute?};
    \end{tikzpicture}
  \end{center}

  \vspace{0.6em}
  For position $i$:
  \begin{enumerate}
    \item compare its query $\mathbf q_i$ with every allowed key $\mathbf k_j$;
    \item turn the comparison scores into attention weights;
    \item use those weights to combine the value vectors $\mathbf v_j$.
  \end{enumerate}

  \begin{alertblock}{Why ``self''-attention?}
    Queries, keys, and values all come from the same sequence $X$.
  \end{alertblock}
\end{frame}

\begin{frame}{Queries, keys, and values are learned projections
  \hfill{\tiny \texttt{Head}: lines 8--10, 17--18, 29}}
  \small
  Three learned matrices give each input row three different roles:
  \[
    Q=XW_Q,\qquad K=XW_K,\qquad V=XW_V.
  \]

  \begin{center}
    \begin{tikzpicture}[
      matrixbox/.style={draw, rounded corners=3pt, minimum height=0.9cm,
        minimum width=2.25cm, align=center},
      >=Latex
    ]
      \node[matrixbox, fill=tensorpurple, draw=tensorpurpleborder] (x)
        {$X$\\$n\times d$};
      \node[matrixbox, fill=tokenorange, draw=tokenorangeborder,
        right=1.2cm of x, yshift=1.25cm] (q) {$Q=XW_Q$\\$n\times d_k$};
      \node[matrixbox, fill=tokenblue, draw=tokenblueborder,
        right=1.2cm of x] (k) {$K=XW_K$\\$n\times d_k$};
      \node[matrixbox, fill=tokengreen, draw=tokengreenborder,
        right=1.2cm of x, yshift=-1.25cm] (v) {$V=XW_V$\\$n\times d_v$};
      \draw[->, thick] (x.east) -- (q.west)
        node[midway, above, sloped] {$W_Q$};
      \draw[->, thick] (x.east) -- (k.west)
        node[midway, above] {$W_K$};
      \draw[->, thick] (x.east) -- (v.west)
        node[midway, below, sloped] {$W_V$};
    \end{tikzpicture}
  \end{center}

  \vspace{-0.35em}
  \begin{columns}[T]
    \begin{column}{0.52\textwidth}
      \begin{block}{Shapes}
        \scriptsize
        \textbf{Weights:}
        $W_Q,W_K\in\mathbb R^{d\times d_k}$,
        $W_V\in\mathbb R^{d\times d_v}$.

        \textbf{Outputs:}
        $Q,K\in\mathbb R^{n\times d_k}$,
        $V\in\mathbb R^{n\times d_v}$.
      \end{block}
    \end{column}
    \begin{column}{0.44\textwidth}
      \begin{alertblock}{How are they learned?}
        \scriptsize
        There are no separate query, key, or value targets. Next-token
        cross-entropy backpropagates through attention and updates
        $W_Q,W_K,W_V$ with all other model weights.
      \end{alertblock}
    \end{column}
  \end{columns}
\end{frame}

\begin{frame}{Step 1: compare one query with every key
  \hfill{\tiny \texttt{Head.forward}: lines 20--22}}
  \small
  Suppose position $i$ contains \texttt{bank}. Its compatibility with an
  allowed position $j$ is the scaled dot product
  \[
    s_{ij}=\frac{\mathbf q_i^{\mathsf T}\mathbf k_j}{\sqrt{d_k}}.
  \]

  \begin{center}
    \begin{tikzpicture}[
      key/.style={draw, rounded corners=3pt, minimum width=2.2cm,
        minimum height=0.72cm, align=center},
      >=Latex
    ]
      \node[key, fill=tokenblue, draw=tokenblueborder] (the)
        {key: \texttt{the}};
      \node[key, fill=tokenorange, draw=tokenorangeborder,
        right=1.0cm of the] (river) {key: \texttt{river}};
      \node[key, fill=tokengreen, draw=tokengreenborder,
        right=1.0cm of river] (bankkey) {key: \texttt{bank}};
      \node[key, fill=tensorpurple, draw=tensorpurpleborder,
        below=1.35cm of river] (bankquery) {query: \texttt{bank}};

      \draw[->, tokenblueborder, line width=0.6pt] (bankquery) --
        node[midway, left] {$s=0.2$} (the);
      \draw[->, tokenorangeborder, line width=2.0pt] (bankquery) --
        node[midway, right] {$s=2.0$} (river);
      \draw[->, tokengreenborder, line width=0.9pt] (bankquery) --
        node[midway, right] {$s=0.6$} (bankkey);
    \end{tikzpicture}
  \end{center}

  \begin{columns}[T]
    \begin{column}{0.48\textwidth}
      \begin{block}{Interpretation}
        A larger score means the query and key are more compatible. A score is
        still an unrestricted number, not a probability.
      \end{block}
    \end{column}
    \begin{column}{0.48\textwidth}
      \begin{block}{Why divide by $\sqrt{d_k}$?}
        Dot products tend to grow as vectors get wider. Scaling keeps scores in
        a range where softmax remains responsive.
      \end{block}
    \end{column}
  \end{columns}
\end{frame}

\begin{frame}{A causal mask prevents looking into the future
  \hfill{\tiny \texttt{Head}: lines 11, 23--25}}
  \small
  During next-token prediction, position $i$ may use only positions
  $j\le i$. Forbidden scores are replaced by $-\infty$ before softmax.

  \begin{columns}[c]
    \begin{column}{0.50\textwidth}
      \[
        S=\frac{QK^{\mathsf T}}{\sqrt{d_k}}+M,
        \qquad
        M_{ij}=\begin{cases}
          0 & j\le i,\\
          -\infty & j>i.
        \end{cases}
      \]
      Because $e^{-\infty}=0$, a masked position receives zero attention
      weight.
    \end{column}
    \begin{column}{0.46\textwidth}
      \centering
      \textbf{Rows attend to columns}

      \vspace{0.4em}
      \begin{tikzpicture}[x=0.78cm,y=0.68cm]
        \foreach \r in {1,...,4} {
          \foreach \c in {1,...,4} {
            \ifnum\c>\r
              \node[draw=gray!55, fill=gray!18, minimum size=0.62cm]
                at (\c,-\r) {$-\infty$};
            \else
              \node[draw=tokenblueborder, fill=tokenblue, minimum size=0.62cm]
                at (\c,-\r) {$s_{\r\c}$};
            \fi
          }
        }
        \foreach \i/\t in {1/the,2/river,3/bank,4/rose} {
          \node[font=\scriptsize, rotate=45, anchor=west] at (\i,-0.35)
            {\texttt{\t}};
          \node[font=\scriptsize, anchor=east] at (0.45,-\i)
            {\texttt{\t}};
        }
      \end{tikzpicture}
    \end{column}
  \end{columns}

  \begin{alertblock}{No information leakage}
    The representation at position $i$ uses tokens $1,\ldots,i$ to predict
    token $i+1$. It cannot depend on token $i+1$ or anything after it.
  \end{alertblock}
\end{frame}

\begin{frame}{Step 2: softmax turns scores into attention weights
  \hfill{\tiny \texttt{Head.forward}: lines 26--27}}
  \small
  For a fixed query position $i$, normalize its allowed scores across $j$:
  \[
    \alpha_{ij}=\frac{e^{s_{ij}}}{\sum_{\ell\le i}e^{s_{i\ell}}},
    \qquad \alpha_{ij}>0,
    \qquad \sum_{j\le i}\alpha_{ij}=1.
  \]

  \begin{center}
    \renewcommand{\arraystretch}{1.25}
    \begin{tabular}{c|c|c|c}
      Key token & Score $s_{ij}$ & $e^{s_{ij}}$ & Weight $\alpha_{ij}$ \\
      \hline
      \texttt{the}   & $0.2$ & $1.22$ & $0.12$ \\
      \texttt{river} & $2.0$ & $7.39$ & $0.71$ \\
      \texttt{bank}  & $0.6$ & $1.82$ & $0.17$ \\
      \texttt{rose}  & $-\infty$ & $0$ & $0$ \\
      \hline
      & & $10.43$ & $1.00$
    \end{tabular}
  \end{center}

  \begin{alertblock}{What the weights say}
    For this illustrative head, \texttt{bank} assigns most of its attention to
    \texttt{river}. Softmax makes the comparison relative: all allowed tokens
    compete for a fixed total weight of $1$.
  \end{alertblock}
\end{frame}

\begin{frame}{Step 3: combine values into contextual information
  \hfill{\tiny \texttt{Head.forward}: lines 29--30}}
  \small
  The weights decide how much of every value vector reaches position $i$:
  \[
    \mathbf z_i=\sum_{j\le i}\alpha_{ij}\mathbf v_j.
  \]

  \begin{center}
    \begin{tikzpicture}[
      value/.style={draw, rounded corners=3pt, minimum width=2.0cm,
        minimum height=0.72cm, fill=tokengreen, draw=tokengreenborder},
      >=Latex
    ]
      \node[value] (the) {$\mathbf v_{\text{the}}$};
      \node[value, right=1.0cm of the] (river) {$\mathbf v_{\text{river}}$};
      \node[value, right=1.0cm of river] (bank) {$\mathbf v_{\text{bank}}$};
      \node[draw=tensorpurpleborder, fill=tensorpurple, rounded corners=3pt,
        below=1.35cm of river, minimum width=2.7cm, minimum height=0.8cm]
        (output) {$\mathbf z_{\text{bank}}$};

      \draw[->, line width=0.5pt, tokenblueborder] (the) --
        node[left, font=\scriptsize] {$0.12$} (output);
      \draw[->, line width=2.2pt, tokenorangeborder] (river) --
        node[right, font=\scriptsize] {$0.71$} (output);
      \draw[->, line width=0.9pt, tokengreenborder] (bank) --
        node[right, font=\scriptsize] {$0.17$} (output);
    \end{tikzpicture}
  \end{center}

  \[
    \mathbf z_{\text{bank}}
    =0.12\mathbf v_{\text{the}}
     +0.71\mathbf v_{\text{river}}
     +0.17\mathbf v_{\text{bank}}.
  \]

  \vspace{-0.7em}
  \begin{alertblock}{The context has entered the vector}
    \footnotesize
    The update for \texttt{bank} now contains substantial information from
    \texttt{river}. Different contexts yield different weights and therefore
    different updates.
  \end{alertblock}
\end{frame}

\begin{frame}{All positions at once: the attention equation
  \hfill{\tiny \texttt{Head.forward}: lines 17--30}}
  \small
  The per-token operations combine into one matrix expression:
  \[
    \boxed{\operatorname{Attention}(Q,K,V)
      =\operatorname{softmax}\!\left(
        \frac{QK^{\mathsf T}}{\sqrt{d_k}}+M
      \right)V}
  \]

  \begin{center}
    \begin{tikzpicture}[
      stage/.style={draw, rounded corners=3pt, minimum width=2.45cm,
        minimum height=1.0cm, align=center},
      >=Latex
    ]
      \node[stage, fill=tokenorange, draw=tokenorangeborder] (scores)
        {scores\\$QK^{\mathsf T}$\\$n\times n$};
      \node[stage, fill=tokenblue, draw=tokenblueborder,
        right=0.65cm of scores] (weights)
        {mask + softmax\\$A$\\$n\times n$};
      \node[stage, fill=tokengreen, draw=tokengreenborder,
        right=0.65cm of weights] (values)
        {mix values\\$AV$\\$n\times d_v$};
      \draw[->, thick] (scores) -- (weights);
      \draw[->, thick] (weights) -- (values);
    \end{tikzpicture}
  \end{center}

  \begin{block}{Reading the shape}
    The $n\times n$ matrix $A$ contains one attention distribution per row.
    Multiplying by $V$ produces one contextual update for every position.
  \end{block}
\end{frame}

\begin{frame}{Multi-head attention learns several relationships
  \hfill{\tiny \texttt{MultiHeadAttention}: lines 34--46}}
  \small
  One attention pattern need not capture every useful relationship. A
  Transformer runs $H$ attention heads with separate learned projections.

  \begin{center}
    \begin{tikzpicture}[
      head/.style={draw, rounded corners=3pt, minimum width=3.05cm,
        minimum height=1.0cm, align=center},
      >=Latex
    ]
      \node[head, fill=tokenblue, draw=tokenblueborder] (h1)
        {head 1\\nearby syntax?};
      \node[head, fill=tokenorange, draw=tokenorangeborder,
        right=0.45cm of h1] (h2) {head 2\\subject relation?};
      \node[head, fill=tokengreen, draw=tokengreenborder,
        right=0.45cm of h2] (h3) {head 3\\semantic clue?};
      \node[draw=tensorpurpleborder, fill=tensorpurple, rounded corners=3pt,
        below=1.05cm of h2, minimum width=4.0cm, minimum height=0.8cm]
        (combine) {concatenate and project with $W_O$};
      \draw[->, thick] (h1) -- (combine);
      \draw[->, thick] (h2) -- (combine);
      \draw[->, thick] (h3) -- (combine);
    \end{tikzpicture}
  \end{center}

  \[
    \operatorname{MultiHead}(X)
      =\operatorname{Concat}(\text{head}_1,\ldots,\text{head}_H)W_O.
  \]

  \begin{alertblock}{Interpret the labels as intuition}
    Heads are not assigned these jobs by hand; training learns their behavior,
    and real heads may mix several kinds of relationships.
  \end{alertblock}
\end{frame}

\begin{frame}{Attention is the communication step in a Transformer block}
  \small
  \begin{center}
    \begin{tikzpicture}[
      stage/.style={draw, rounded corners=3pt, minimum width=2.15cm,
        minimum height=0.78cm, align=center},
      >=Latex
    ]
      \node[stage, fill=tokenblue, draw=tokenblueborder] (input) {$X$};
      \node[stage, fill=tokenorange, draw=tokenorangeborder,
        right=0.38cm of input] (attention) {causal\\self-attention};
      \node[stage, fill=tensorpurple, draw=tensorpurpleborder,
        right=0.38cm of attention] (res1) {add residual\\+ normalize};
      \node[stage, fill=tokengreen, draw=tokengreenborder,
        right=0.38cm of res1] (mlp) {per-token\\MLP};
      \node[stage, fill=tensorpurple, draw=tensorpurpleborder,
        right=0.38cm of mlp] (res2) {add residual\\+ normalize};
      \draw[->, thick] (input) -- (attention);
      \draw[->, thick] (attention) -- (res1);
      \draw[->, thick] (res1) -- (mlp);
      \draw[->, thick] (mlp) -- (res2);
    \end{tikzpicture}
  \end{center}

  \vspace{0.8em}
  \begin{columns}[T]
    \begin{column}{0.48\textwidth}
      \begin{block}{Across positions}
        Self-attention moves information between tokens. The causal mask
        preserves next-token prediction.
      \end{block}
    \end{column}
    \begin{column}{0.48\textwidth}
      \begin{block}{Within each position}
        The MLP transforms each token vector independently. Residual paths help
        preserve and refine existing information.
      \end{block}
    \end{column}
  \end{columns}

  \begin{alertblock}{Depth builds richer context}
    Stacking many blocks lets information travel and be transformed repeatedly,
    producing the contextual representations used to predict the next token.
  \end{alertblock}
\end{frame}

\begin{frame}{Inside the per-token MLP: factual associations}
  \small
  \begin{center}
    \begin{tikzpicture}[
      stage/.style={draw, rounded corners=3pt, minimum height=1.0cm,
        minimum width=2.55cm, align=center},
      >=Latex
    ]
      \node[stage, fill=tokenblue, draw=tokenblueborder] (context)
        {one token's context\\``Michael Jordan plays \ldots{}''};
      \node[stage, fill=tokenorange, draw=tokenorangeborder,
        right=0.55cm of context] (key)
        {\(W_{\mathrm{up}}\)\\pattern match \(a_i\)};
      \node[stage, fill=tensorpurple, draw=tensorpurpleborder,
        right=0.55cm of key] (value)
        {\(W_{\mathrm{down}}\)\\add value \(\mathbf v_i\)};
      \node[stage, fill=tokengreen, draw=tokengreenborder,
        right=0.55cm of value] (prediction)
        {higher score for\\\texttt{basketball}};
      \draw[->, thick] (context) -- (key);
      \draw[->, thick] (key) -- (value);
      \draw[->, thick] (value) -- (prediction);
    \end{tikzpicture}
  \end{center}

  \vspace{-0.5em}
  \[
    \operatorname{MLP}(\mathbf x)
      =\sigma(\mathbf xW_{\mathrm{up}})W_{\mathrm{down}}
      =\sum_i a_i\mathbf v_i,
    \qquad
    a_i=\sigma(\mathbf x\mathbf k_i).
  \]
  \vspace{-1.0em}
  \begin{center}
    \scriptsize
    \(\mathbf k_i\) is column \(i\) of \(W_{\mathrm{up}}\);
    \(\mathbf v_i\) is row \(i\) of \(W_{\mathrm{down}}\).
  \end{center}

  \vspace{-0.4em}
  \begin{columns}[T,onlytextwidth]
    \begin{column}{0.49\textwidth}
      \begin{block}{Geva et al. (2021)}
        \scriptsize
        \href{https://arxiv.org/abs/2012.14913}{\emph{Transformer Feed-Forward
        Layers Are Key-Value Memories}.}
        They found that \(W_{\mathrm{up}}\) directions respond to interpretable
        textual patterns, while corresponding \(W_{\mathrm{down}}\) values favor
        plausible output tokens.
      \end{block}
    \end{column}
    \begin{column}{0.49\textwidth}
      \begin{block}{Meng et al. (2022): ROME}
        \scriptsize
        \href{https://arxiv.org/abs/2202.05262}{\emph{Locating and Editing
        Factual Associations in GPT}.}
        Causal tracing localized factual predictions to middle-layer MLPs;
        rank-one changes to their feed-forward weights could update particular
        factual associations.
      \end{block}
    \end{column}
  \end{columns}

  \vspace{-0.4em}
  \begin{alertblock}{Important nuance}
    \footnotesize
    ``Michael Jordan plays basketball'' is illustrative. Evidence points to
    distributed, compositional storage---not one complete sentence in one neuron.
  \end{alertblock}
\end{frame}

\begin{frame}{GPT-3: sizes of the attention projection matrices}
  \small
  GPT-3 uses \(H=96\) attention heads in each of \(N=96\) layers, with

  \[
    d_q=d_k=d_v=128,
    \qquad
    d_{\mathrm{embed}}=12{,}288.
  \]

  \vspace{-0.8em}
  \begin{block}{From separate heads to one combined projection}
    \centering
    \(H d_q=H d_k=H d_v=96\cdot128=12{,}288=d_{\mathrm{embed}}\).
    Thus each combined projection matrix is square.
  \end{block}

  \begin{center}
    \scriptsize
    \renewcommand{\arraystretch}{1.3}
    \begin{tabular}{@{}lccc@{}}
      \toprule
      Matrix & Shape in one layer & Parameters per layer & Across 96 layers \\
      \midrule
      \(W_Q\) & \(12{,}288\times12{,}288\) & \(150{,}994{,}944\) & \(14{,}495{,}514{,}624\) \\
      \(W_K\) & \(12{,}288\times12{,}288\) & \(150{,}994{,}944\) & \(14{,}495{,}514{,}624\) \\
      \(W_V\) & \(12{,}288\times12{,}288\) & \(150{,}994{,}944\) & \(14{,}495{,}514{,}624\) \\
      \(W_O\) & \(12{,}288\times12{,}288\) & \(150{,}994{,}944\) & \(14{,}495{,}514{,}624\) \\
      \bottomrule
    \end{tabular}
  \end{center}

  \vspace{-0.4em}
  \begin{alertblock}{Attention projections across the whole model}
    \centering
    \(4\times14{,}495{,}514{,}624
      =57{,}982{,}058{,}496\) learned weights, excluding biases.
  \end{alertblock}
\end{frame}

\begin{frame}{GPT-3: the MLP up and down projections}
  \small
  Attention moves information between tokens. The \textbf{MLP} then transforms
  each token vector independently, using the same two matrices at every position.

  \begin{center}
    \begin{tikzpicture}[
      stage/.style={draw, rounded corners=3pt, minimum height=0.95cm,
        align=center},
      >=Latex
    ]
      \node[stage, minimum width=1.55cm, fill=tokenblue,
        draw=tokenblueborder] (input)
        {token vector\\\(d=12{,}288\)};
      \node[stage, minimum width=2.0cm, fill=tokenorange,
        draw=tokenorangeborder, right=0.35cm of input] (up)
        {up-project\\with \(W_{\mathrm{up}}\)};
      \node[stage, minimum width=1.7cm, fill=tokengreen,
        draw=tokengreenborder, right=0.35cm of up] (wide)
        {wide vector\\\(d_{\mathrm{ff}}=49{,}152\)};
      \node[stage, minimum width=1.25cm, fill=tensorpurple,
        draw=tensorpurpleborder, right=0.35cm of wide] (activation)
        {nonlinear\\activation};
      \node[stage, minimum width=2.0cm, fill=tokenorange,
        draw=tokenorangeborder, right=0.35cm of activation] (down)
        {down-project\\with \(W_{\mathrm{down}}\)};
      \node[stage, minimum width=1.55cm, fill=tokenblue,
        draw=tokenblueborder, right=0.35cm of down] (output)
        {MLP update\\\(d=12{,}288\)};
      \draw[->, thick] (input) -- (up);
      \draw[->, thick] (up) -- (wide);
      \draw[->, thick] (wide) -- (activation);
      \draw[->, thick] (activation) -- (down);
      \draw[->, thick] (down) -- (output);
    \end{tikzpicture}
  \end{center}

  \vspace{-0.5em}
  \begin{center}
    \scriptsize
    \renewcommand{\arraystretch}{1.3}
    \begin{tabular}{@{}lccc@{}}
      \toprule
      Matrix & Shape in one layer & Parameters per layer & Across 96 layers \\
      \midrule
      \(W_{\mathrm{up}}\) & \(12{,}288\times49{,}152\)
        & \(603{,}979{,}776\) & \(57{,}982{,}058{,}496\) \\
      \(W_{\mathrm{down}}\) & \(49{,}152\times12{,}288\)
        & \(603{,}979{,}776\) & \(57{,}982{,}058{,}496\) \\
      \bottomrule
    \end{tabular}
  \end{center}

  \vspace{-0.4em}
  \begin{alertblock}{Most GPT-3 parameters are in its MLPs}
    \centering
    Together: \(2\times57{,}982{,}058{,}496
      =115{,}964{,}116{,}992\) weights---about \(66\%\) of GPT-3.
  \end{alertblock}
\end{frame}

\begin{frame}{Model scale: GPT-3 175B (2020) vs. Gemma 4 31B (2026)}
  \small
  Both are Transformer models, but Gemma 4 reaches a much smaller parameter
  budget through a narrower, more modern architecture.

  \begin{center}
    \begin{tikzpicture}[font=\small]
      \node[anchor=east] at (0,0.55) {GPT-3};
      \fill[tokenblueborder] (0,0.25) rectangle (10.5,0.85);
      \node[anchor=west, text=white, font=\bfseries]
        at (0.15,0.55) {175.18B parameters};

      \node[anchor=east] at (0,-0.35) {Gemma 4};
      \fill[tokenorangeborder] (0,-0.65) rectangle (1.84,-0.05);
      \node[anchor=west, font=\bfseries]
        at (2.02,-0.35) {30.7B parameters};
    \end{tikzpicture}
  \end{center}

  \vspace{-0.7em}
  \begin{center}
    \scriptsize
    \renewcommand{\arraystretch}{1.22}
    \begin{tabular}{@{}lrrrr@{}}
      \toprule
      Parameter family & GPT-3 175B & Share &
        \href{https://chatgpt.com/share/6a56a6e4-0c14-83ed-849e-93f362e5f1c3}{Gemma 4 estimate} & Share \\
      \midrule
      Embedding and output & \(1.24\)B & \(0.7\%\) & \(\approx1.41\)B & \(\approx4.6\%\) \\
      Attention: \(Q,K,V,O\) & \(57.98\)B & \(33.1\%\) & \(\approx6.10\)B & \(\approx19.8\%\) \\
      MLP: up and down projections & \(115.96\)B & \(66.2\%\) & \(\approx13.88\)B & \(\approx45.2\%\) \\
      Vision encoder and other weights & \(\approx0\)B & \(\approx0\%\) & \(\approx9.35\)B & \(\approx30.4\%\) \\
      \midrule
      \textbf{Total} & \textbf{175.18B} & \textbf{100\%} &
        \href{https://ai.google.dev/gemma/docs/core/model_card_4}{\textbf{30.7B}} & \textbf{\(\approx100\%\)} \\
      \bottomrule
    \end{tabular}
  \end{center}

  \vspace{-0.6em}
  \begin{alertblock}{About \(5.7\times\) fewer parameters}
    \footnotesize
    Gemma 4 uses 60 layers, a smaller embedding dimension, grouped-query and
    hybrid local/global attention, plus a vision encoder. Size alone is not a
    measure of capability; architecture, data, and training matter enormously.
  \end{alertblock}

\end{frame}
